\subsection{Créditos en estudios avanzados}
Acreditas módulos o unidades de un diplomado superior, especialidad o maestría afín a tu licenciatura. (Art. 22, REP).

\subsubsection*{Requisitos}
\begin{itemize}
  \item Presentar la solicitud y obtener la autorización de la Subdirección Académica antes de inscribirte a los estudios.
  \item Que la Coordinación de Docencia dictamine la afinidad con tu licenciatura.
  \item Contar con carta de aceptación y con el plan de estudios oficial aprobado por los Consejos.
  \item Certificar el 100 \% de los créditos (diplomado superior o especialidad) o el 75 \% (maestría).
\end{itemize}
\subsubsection*{Etapa 1. Tus datos}
Captura una sola vez tus datos; con ellos se llenan automáticamente todos los formatos del trámite.

Datos que se capturan:
\begin{itemize}
  \item Nombre completo (con acentos, tal como aparecerá en el título)
  \item Número de cuenta
  \item Licenciatura
  \item Sistema
  \item Plan de estudios (clave y versión)
  \item Escuela incorporada a la UAEM donde cursaste la licenciatura
  \item Año de ingreso
  \item Año de egreso
  \item Teléfono
  \item Correo institucional
  \item Correo personal
  \item Ocupación actual
  \item Fecha de nacimiento
  \item Mi nombre o apellidos llevan acentos
\end{itemize}
\subsubsection*{Etapa 2. Solicitud de la opción}
Presenta la solicitud antes de inscribirte a los estudios avanzados. La Coordinación de Docencia dictamina la afinidad y la Subdirección Académica autoriza.

Datos que se capturan:
\begin{itemize}
  \item La solicitud se presenta por
  \item Estudios avanzados que cursarás (programa, institución y, si aplica, RVOE)
\end{itemize}

La plataforma llena los formatos: creditos y los firma por ti.

Documentos (PDF):
\begin{longtable}{p{0.55\linewidth}p{0.35\linewidth}}
\toprule
\textbf{Documento} & \textbf{Observaciones}\\
\midrule
\endhead
Solicitud firmada por el pasante & La plataforma lo genera. Debe estar fechada antes de iniciar los estudios avanzados.\\
Certificado de estudios de licenciatura & \\
Certificado de servicio social & \\
Plan de estudios oficial de los estudios avanzados & Debe avalar cursos, créditos, actividades con valor curricular, plazos mínimos, perfiles y contenidos programáticos.\\
Registro del programa ante la Dirección General de Profesiones (o validez oficial) & Solo para estudios que no se realicen en la UAEMéx.\\
\bottomrule
\end{longtable}

\emph{Turno del departamento:} Dictamen de afinidad y autorización. La Coordinación de Docencia dictamina la afinidad y la Subdirección Académica autoriza.
\subsubsection*{Etapa 3. Estudios avanzados cursados}
Cuando concluyas, sube la documentación que acredita el 100 \% de los créditos (diplomado superior o especialidad) o el 75 \% (maestría), cubiertos en los plazos mínimos.

Datos que se capturan:
\begin{itemize}
  \item Programa y institución donde cursaste
\end{itemize}

Documentos (PDF):
\begin{longtable}{p{0.55\linewidth}p{0.35\linewidth}}
\toprule
\textbf{Documento} & \textbf{Observaciones}\\
\midrule
\endhead
Carta de aceptación a los estudios avanzados & Con fecha posterior a la carta de afinidad.\\
Dictamen de equivalencia de la Secretaría de Investigación y Estudios Avanzados & \\
Certificado de estudios y constancia de calificaciones de créditos cubiertos & Especifica que se cubrieron en los plazos mínimos.\\
\bottomrule
\end{longtable}
\subsubsection*{Etapa 4. Documentación para la toma de protesta}
Sube cada documento en PDF, legible y a color. Cuando todo esté validado, Apoyo a Titulación te citará para entregar originales y copias y fijar la fecha.

La plataforma llena los formatos: datos\_titulo, requisitos y los firma por ti.

Documentos (PDF):
\begin{longtable}{p{0.55\linewidth}p{0.35\linewidth}}
\toprule
\textbf{Documento} & \textbf{Observaciones}\\
\midrule
\endhead
Certificado de estudios de bachillerato & \\
Constancia de autenticación de bachillerato & Se tramita en la escuela donde cursaste el bachillerato. Solo si tus estudios de bachillerato no fueron en instituciones dependientes o incorporadas a la UAEM.\\
Acta de nacimiento & Copia certificada vigente (del año en curso).\\
CURP & Debe mostrar la leyenda «Validada por el Registro Civil».\\
Constancia o certificado de aprobación del idioma inglés & Solo planes 98 en adelante, si en tu certificado no aparecen Inglés C1 e Inglés C2.\\
Constancia de no adeudo a la Biblioteca de la Facultad & El recibo de pago se tramita en Control Escolar de la Facultad; la constancia, en la Biblioteca.\\
Constancia de no adeudo a la Biblioteca Central & \\
Contra-recibo de la donación de libros a la Biblioteca de la Facultad & \\
Constancia de no adeudo a la Dirección de Control Escolar de la UAEM & Se solicita en el Portal de alumnos: controlescolar.uaemex.mx.\\
Recibo de no adeudo a la UAEM & Se solicita en el Portal de alumnos de Control Escolar.\\
Formato de datos personales (para el título) & La plataforma lo genera. La plataforma lo llena con tus datos; lo firmas electrónicamente.\\
Comprobante del sistema de Seguimiento de Egresados & Se obtiene en alumniversitario.uaemex.mx/1dae.php.\\
Dictamen de autorización de los HH. Consejos & La plataforma lo genera. Se solicita en la Subdirección Académica con el formato del Art. 5. Solo si excediste el plazo de 10 años para presentar la evaluación profesional (Art. 5, REP).\\
Recibo de pago: derechos de evaluación profesional (\$230) & Se solicita en Control Escolar de la Facultad.\\
Recibo de pago: expedición de título (\$1,187) & Se solicita en Control Escolar de la Facultad.\\
\bottomrule
\end{longtable}
\subsubsection*{Etapa 5. Fecha, hora y lugar de la toma de protesta}
Apoyo a Titulación asigna el jurado y te comunica la fecha. Preséntate puntualmente.

\emph{Turno del departamento:} Fecha y lugar de la toma de protesta. Titulación fija la fecha. En esta opción la comisión es siempre el Director (presidente) y el Subdirector Académico (secretario); el acta oficial se genera en otro sistema y se firma en papel.
\subsubsection*{Etapa 6. Resultado de la evaluación}
Al concluir, el departamento registra el resultado. Después podrás consultar cómo descargar tu título y cédula profesional.

\emph{Turno del departamento:} Registrar resultado. 

\subsubsection*{Seguimiento y conclusión}
\subsubsection*{Pantallas del proceso}
\captura{creditos_pasante_01_opciones}{Opciones de titulación en «Licenciatura».}
\captura{creditos_pasante_02_requisitos}{Requisitos de la opción y aceptación para iniciar.}
\captura{creditos_pasante_03_iniciar}{Aceptación de condiciones para iniciar el trámite.}
\captura{creditos_pasante_04_datos__1}{Captura de datos del pasante y apertura del trámite. (parte 1 de 2)}
\captura{creditos_pasante_04_datos__2}{Captura de datos del pasante y apertura del trámite. (parte 2 de 2)}
\captura{creditos_pasante_05_ronda0__1}{Etapa en curso del trámite. (parte 1 de 3)}
\captura{creditos_pasante_05_ronda0__2}{Etapa en curso del trámite. (parte 2 de 3)}
\captura{creditos_pasante_05_ronda0__3}{Etapa en curso del trámite. (parte 3 de 3)}
\captura{creditos_pasante_06_concluido__1}{Trámite concluido. (parte 1 de 3)}
\captura{creditos_pasante_06_concluido__2}{Trámite concluido. (parte 2 de 3)}
\captura{creditos_pasante_06_concluido__3}{Trámite concluido. (parte 3 de 3)}
\captura{creditos_pasante_07_ruta__1}{Ruta del trámite con los pasos hechos. (parte 1 de 2)}
\captura{creditos_pasante_07_ruta__2}{Ruta del trámite con los pasos hechos. (parte 2 de 2)}
\captura{creditos_pasante_08_estado}{Estado del trámite.}
